\section{Related Work}

\paragraph{Embedding Distillation Beyond Coordinate Matching.}
Knowledge distillation transfers behavior from a large teacher to a smaller student~\citep{hinton2015distilling}. In pretrained language models, DistilBERT, TinyBERT, and MiniLM compress transformer encoders by matching logits, hidden states, attention distributions, or self-attention relations~\citep{sanh2019distilbert,jiao2019tinybert,wang2020minilm,xu2024survey}. These methods provide effective compression recipes, but many rely on teacher internals, token-level correspondences, or output distributions that are not always available for black-box embedding teachers.

Embedding-specific distillation usually supervises final representations or task-derived relations. Sentence embedding methods such as DistilCSE transfer contrastive behavior and pairwise preferences~\citep{gao2021distilcse,xu2023distillcse}. Gecko uses LLM-generated and LLM-relabelled retrieval data, while Jasper and Stella combine multi-stage embedding distillation with Matryoshka Representation Learning~\citep{lee2024gecko,zhang2024jasperstella}. LEAF supports asymmetric retrieval by aligning compact students to teacher representations, and Jina Embeddings v5 combines task-targeted distillation with contrastive training for compact multilingual embeddings~\citep{vujanic2025leaf,akram2026jinaembeddingsv5}. These approaches establish distillation as a central tool for embedding model development. MDD differs in the object being transferred: instead of matching individual coordinates, generated pairs, or mini-batch similarities, it matches multi-scale diffusion profiles induced by the teacher geometry.

\paragraph{Diffusion Geometry and Manifold Learning.}
Graph-based manifold learning studies how local neighborhoods reveal low-dimensional structure in high-dimensional data. Laplacian Eigenmaps preserve local neighborhoods through graph Laplacian eigenvectors~\citep{belkin2003laplacian}, while Diffusion Maps use Markov diffusion on a neighborhood graph to expose multi-scale connectivity and intrinsic geometry~\citep{coifman2006diffusionmaps}. In this view, the important object is not an arbitrary coordinate system but the diffusion operator, its neighborhood expansion behavior, and the low-frequency modes of the graph.

MDD adapts this perspective to text embedding distillation. The teacher embeddings define finite support samples from a semantic manifold, and a sparse mutual $k$NN graph provides a discrete approximation of that manifold. Powers of the Markov transition matrix define diffusion profiles that describe how semantic mass spreads from an anchor to local neighbors, broader regions, and higher-order communities. The spectral anchor complements this local diffusion supervision by transferring low-frequency Laplacian coordinates. Thus, the student is trained to preserve teacher-induced diffusion geometry rather than only reproduce teacher coordinates.

\paragraph{Geometry-Preserving Representation Transfer.}
Relational knowledge distillation argues that students should preserve relations among examples rather than only individual outputs~\citep{park2019rkd}. In embedding training, this often becomes mini-batch similarity matching, where teacher and student pairwise similarity matrices are compared over the examples sampled in a batch. This improves over pure coordinate regression, but the supervision remains local to the sampled batch and does not explicitly model multi-hop semantic paths or broader community structure.

Other representation alignment tools, including CKA-style similarity analysis, optimal-transport alignment, EMO-style objectives, and teacher-anchored layer alignment, provide useful ways to compare or transfer representations~\citep{kornblith2019similarity,alqahtani2021ot,truong2025emo,dao2026talas}. These methods are valuable when intermediate layers or alignment pairs are available, but they generally optimize local representational agreement rather than a coordinate-free manifold object. MDD keeps the architecture-agnostic advantage of final-embedding distillation while replacing batch-local relational snapshots with cached diffusion targets sampled from the teacher-induced support graph.

\paragraph{Text Embedding Evaluation as Probes.}
Text embedding models are commonly evaluated through retrieval, classification, clustering, semantic textual similarity, and transfer benchmarks. Sentence-BERT and SimCSE made sentence-level similarity evaluation central to embedding research~\citep{reimers2019sentencebert,gao2021simcse}, while BEIR and MTEB broadened evaluation toward multi-task retrieval and representation assessment~\citep{thakur2021beir,muennighoff2023mteb}. Continued maintenance of MTEB reflects the need for reproducible embedding evaluation as models and datasets evolve~\citep{chung2025maintainingmteb}. Recent large embedding systems such as BGE-M3, Qwen3 Embedding, and Dewey provide strong teacher candidates and broad deployment motivation~\citep{chen2024bgem3,zhang2025qwen3embedding,zhang2025dewey}.

Our experiments use classification, pair classification, and semantic textual similarity as frozen-representation probes rather than as claims of exhaustive downstream coverage. Classification probes test whether class-level semantic regions remain linearly separable; pair classification probes test local semantic discrimination; and STS probes test continuous similarity calibration. This evaluation matches our central claim: preserving teacher-induced diffusion geometry should yield compact embeddings that transfer across complementary aspects of semantic representation quality.
